\casechapter{All Minimax Theorems Are One}{Linear programming duality}{Duality}{ch:lpduality}
\begin{refsection}[chapters/05-lpduality.bib]

\newcommand{\lpone}{\mathbf{1}}
\newcommand{\lpcone}{\operatorname{cone}}
\newcommand{\lpval}{\operatorname{val}}
\newcommand{\lpcost}{\operatorname{cost}}
\newcommand{\lperr}{\operatorname{err}}
\newcommand{\lpQ}{\mathcal{Q}}

\begin{chaptersummary}
Between 1927 and 1956 a series of theorems appeared with one shape: the largest matching in a bipartite graph equals the smallest vertex cover, the largest flow equals the smallest cut, the largest antichain equals the fewest chains, and in a zero-sum game the best guaranteed gain equals the best guaranteed loss, each with its own proof. The move is to write both sides as linear programs: every easy half becomes weak duality, every hard half becomes the duality theorem of linear programming, proved once from a separating hyperplane, and integrality becomes a determinant condition on the constraint matrix. The zoo dissolves into one theorem, which then reaches randomized algorithms, boosting, and hardness amplification.
\end{chaptersummary}

\section{Before: a zoo of min--max theorems}\label{sec:lpduality:before}

On 26 March 1931, D\'enes K\H{o}nig told the Budapest Mathematical and Physical Society that in a bipartite graph the largest number of pairwise disjoint edges equals the smallest number of vertices touching every edge.\sidenote{The date is recorded in \textcite{Schrijver2005}, a history of combinatorial optimization up to 1960 and the source of several historical details in this chapter.} Between 1927 and 1956 the same shape appeared again and again (\cref{tab:lpduality:zoo}): a maximum over one kind of object equals a minimum over another kind.\sidenote{The max-flow problem came to Ford and Fulkerson from T.~E.~Harris at the RAND Corporation. A 1955 RAND report by Harris and F.~S.~Ross on the railways of the Soviet Union and Eastern Europe, declassified in 1999 at Schrijver's request, computed a flow and a cut (``the bottleneck'') of equal capacity; the interest was in the cut \parencite{Schrijver2005}.} Wherever a maximum meets a minimum, the inequality ``maximum $\le$ minimum'' is a one-line count, and the reverse inequality, the existence of an object that certifies optimality, needed a new idea.

\begin{table}[t]
\centering\small
\begin{tabularx}{\linewidth}{@{}l l X@{}}
\toprule
Theorem & Year & The largest \dots{} equals the smallest \dots\\
\midrule
Menger & 1927 & number of disjoint $P$--$Q$ paths; number of vertices meeting all $P$--$Q$ paths\\
von Neumann & 1928 & guaranteed expected gain; guaranteed expected loss\\
K\H{o}nig & 1931 & size of a matching; size of a vertex cover\\
Egerv\'ary & 1931 & weight of an assignment; total of a covering labeling\\
Birkhoff & 1946 & (doubly stochastic matrices are mixtures of permutation matrices)\\
Dilworth & 1950 & size of an antichain; number of chains covering the order\\
Max-flow min-cut & 1956 & value of a flow; capacity of a cut\\
\bottomrule
\end{tabularx}
\caption{The zoo \parencite{Menger1927,vonNeumann1928,Konig1931,Egervary1931,Birkhoff1946,Dilworth1950}. The max-flow min-cut theorem was found independently by \textcite{FordFulkerson1956} and by \textcite{EliasFeinsteinShannon1956}.}\label{tab:lpduality:zoo}
\end{table}

\subsection{K\H{o}nig's theorem by augmenting paths}

\begin{definition}[Matchings and covers]\label{def:lpduality:matching}
A finite graph $G=(V,E)$ is \term{bipartite} with sides $L$ and $R$ if $V$ is the disjoint union of $L$ and $R$ and every edge joins a vertex of $L$ to a vertex of $R$. A \term{matching} is a set of pairwise disjoint edges, and a \term{vertex cover} is a set of vertices that meets every edge. We write $\nu(G)$ for the largest size of a matching and $\tau(G)$ for the smallest size of a vertex cover.
\end{definition}

The inequality $\nu(G)\le\tau(G)$ holds in every graph: a vertex cover contains an endpoint of each edge of a matching, and these endpoints are distinct because the edges are disjoint.\index{K\H{o}nig's theorem}

\begin{theorem}[\textcite{Konig1931}]\label{thm:lpduality:konig}
Every bipartite graph $G$ satisfies $\nu(G)=\tau(G)$.
\end{theorem}

The classical proof grows a cover out of a maximum matching $M$.\sidenote{Egerv\'ary's weighted proof is essentially an algorithm. In 1955 Harold Kuhn sharpened it into the Hungarian method for the assignment problem, named for its Hungarian sources \parencite{Schrijver2005}.} A path is \term{$M$-alternating} if its edges lie alternately outside and inside $M$, and \term{$M$-augmenting} if, in addition, its two endpoints are distinct and touched by no edge of $M$. Exchanging the edges inside and outside $M$ along an augmenting path produces a matching with one more edge, so a maximum matching has no augmenting path.

\begin{proof}
Let $M$ be a maximum matching, let $U$ be the set of vertices of $L$ touched by no edge of $M$, and let $Z$ be the set of vertices reachable from $U$ by $M$-alternating paths, including paths of length zero, so that $U\subseteq Z$. The plan is to show that $C=(L\setminus Z)\cup(R\cap Z)$ is a vertex cover with $|C|\le|M|$.

\emph{Two facts.} An alternating path that starts in $U$ leaves $L$ along edges outside $M$, because its first edge is outside $M$, and it returns to $L$ along edges of $M$. Hence every vertex of $R\cap Z$ is touched by $M$, since otherwise an alternating path from $U$ to it would be augmenting. And if $\ell r\in M$ and $r\in Z$, then $\ell\in Z$: a path from $U$ to $r$ arrives along an edge outside $M$, so it either contains $\ell$ already or extends by the edge $r\ell$.

\emph{The set $C$ is a cover.} Let $\ell r$ be an edge with $\ell\in L$. If $\ell\notin Z$, then $\ell\in C$. If $\ell\in Z$ and $\ell r\notin M$, then a path from $U$ to $\ell$ has length zero or arrives along an edge of $M$, so it contains $r$ or extends by the edge $\ell r$; thus $r\in Z$. If $\ell\in Z$ and $\ell r\in M$, then $\ell\notin U$, and a path from $U$ to $\ell$ arrives along the only edge of $M$ at $\ell$, namely $\ell r$; again $r\in Z$. In every case $\ell\in C$ or $r\in R\cap Z\subseteq C$.

\emph{The size of $C$.} Every vertex of $L\setminus Z$ is touched by $M$ because $U\subseteq Z$, and every vertex of $R\cap Z$ is touched by $M$ by the first fact. No edge of $M$ has both endpoints in $C$, by the second fact. So the map that sends each vertex of $C$ to the edge of $M$ touching it is injective, and $\tau(G)\le|C|\le|M|=\nu(G)$.
\end{proof}

\begin{figure}[t]
\centering
\begin{tikzpicture}[x=1cm,y=0.62cm,
  v/.style={circle,draw=sea,line width=0.6pt,fill=white,inner sep=0pt,minimum size=5.5pt},
  c/.style={circle,draw=seadeep,fill=seadeep,inner sep=0pt,minimum size=5.5pt}]
  \foreach \i/\y in {1/3,2/2,3/1,4/0} {\coordinate (a\i) at (0,\y); \coordinate (b\i) at (2.6,\y);}
  \draw[muted,line width=0.5pt] (a1)--(b1) (a1)--(b2) (a3)--(b2) (a3)--(b4);
  \draw[sea,line width=2pt] (a2)--(b1) (a3)--(b3) (a4)--(b2);
  \node[v] at (a1) {}; \node[v] at (a2) {}; \node[c] at (a3) {}; \node[v] at (a4) {};
  \node[c] at (b1) {}; \node[c] at (b2) {}; \node[v] at (b3) {}; \node[v] at (b4) {};
  \foreach \i in {1,2,3,4} {\node[left=3pt,font=\small] at (a\i) {$\ell_\i$}; \node[right=3pt,font=\small] at (b\i) {$r_\i$};}
  \node[font=\small,text=muted] at (0,3.9) {$L$};
  \node[font=\small,text=muted] at (2.6,3.9) {$R$};
\end{tikzpicture}
\caption{Two certificates. The three thick edges form a matching, and the three filled vertices form a vertex cover. Each certifies that the other is optimal.}\label{fig:lpduality:konig}
\end{figure}

The proof is an algorithm, which is its virtue. But bipartiteness enters invisibly, where alternating paths return to $L$ only along edges of $M$, and the cover comes from a construction made for this one theorem.

\subsection{Doubly stochastic matrices}

A real square matrix is \term{doubly stochastic} if its entries are nonnegative and each of its rows and columns sums to $1$. Convex combinations of permutation matrices are doubly stochastic, and the converse, often named after Birkhoff and von Neumann, belongs to the zoo because it makes the assignment problem a linear program with an integral optimum, as Dantzig observed \parencite{Schrijver2005}.\index{Birkhoff--von Neumann theorem}

\begin{theorem}[\textcite{Birkhoff1946}]\label{thm:lpduality:bvn}
Every doubly stochastic matrix is a convex combination of permutation matrices.
\end{theorem}

\begin{proof}
The plan is to find a permutation matrix under the support of a doubly stochastic $n\times n$ matrix $X$, subtract as much of it as possible, and induct on the number of nonzero entries.

Let $G$ be the bipartite graph whose sides are the rows and the columns of $X$, with row $i$ joined to column $j$ when $X_{ij}>0$. A vertex cover of $G$ is a set of rows and columns containing all nonzero entries, so the total sum $n$ of the entries is at most the number of lines in the cover, each line summing to $1$. Hence $\tau(G)\ge n$, and \cref{thm:lpduality:konig} gives a matching of size $n$: a permutation $\sigma$ with $X_{i\sigma(i)}>0$ for every $i$. Let $P$ be its permutation matrix and $\lambda=\min_i X_{i\sigma(i)}\in(0,1]$. If $\lambda=1$, then each row of $X$ has all its mass at the entry $\sigma(i)$, and $X=P$. If $\lambda<1$, then $X'=(X-\lambda P)/(1-\lambda)$ is doubly stochastic, by the choice of $\lambda$, and has fewer nonzero entries than $X$. By induction $X'$ is a convex combination of permutation matrices, and so is $X=\lambda P+(1-\lambda)X'$.
\end{proof}

\subsection{Games, and a gap in Menger's proof}

One member of the family comes from game theory. A real $m\times n$ matrix $A$ defines a game in which a row player chooses a row $i$, a column player simultaneously chooses a column $j$, and the column player pays $A_{ij}$ to the row player. Write $\Delta_k$ for the set of probability vectors in $\R^k$; mixed strategies $p\in\Delta_m$ and $q\in\Delta_n$ give the expected payment $p^\top Aq$.\index{minimax theorem}

\begin{theorem}[\textcite{vonNeumann1928}]\label{thm:lpduality:minimax}
For every real $m\times n$ matrix $A$,
\[
\max_{p\in\Delta_m}\,\min_{q\in\Delta_n} p^\top Aq=\min_{q\in\Delta_n}\,\max_{p\in\Delta_m} p^\top Aq .
\]
\end{theorem}

The left side is what the row player can guarantee by announcing a mixed strategy, the right side is what the column player can guarantee, and the inequality $\le$ is again the easy half. Von Neumann's first proofs drew on topology and fixed-point arguments; the first elementary proof, by Jean Ville in 1938, used convexity \parencite{Ville1938,Kjeldsen2001}.

\begin{aside}[Menger's gap]
\textcite{Menger1927} proved that the largest number of vertex-disjoint paths from a vertex set $P$ to a vertex set $Q$ equals the smallest number of vertices meeting all such paths. His proof was an induction on the number of edges, and he treated the base case, in which every vertex lies in $P\cup Q$, as easy. It is not, as K\H{o}nig later noticed \parencite{Schrijver2005}: there the minimal paths are the edges between $P\setminus Q$ and $Q\setminus P$, together with the vertices of $P\cap Q$, so the base case is K\H{o}nig's theorem. The min--max content lived in the case the proof skipped.
\end{aside}

\subsection{What the old proofs could not see}

Each classical proof explains its own theorem, and none explains the family. Why the same shape? Where do the integers come from? Which hypothesis does the work, when von Neumann's theorem holds for every real matrix while K\H{o}nig's can fail outside bipartite graphs (a triangle has $\nu=1$ and $\tau=2$)? And why must each generalization, to weights or capacities, pay for a new proof?

\section{The move: write both sides as linear programs}\label{sec:lpduality:move}

Look at the easy halves again. Each edge of a matching is charged to a cover vertex at one of its ends; each unit of flow is charged to an arc leaving the cut; each element of an antichain is charged to the chain containing it. Every easy half is a weighted count in which objects of one kind pay for objects of the other. The move is to take that count seriously, as an inequality between linear functions, and to let the objects become fractional, so that both sides live in vector spaces where convexity is available.\index{linear programming}

\begin{definition}[Linear programs in duality]\label{def:lpduality:pair}
Let $A\in\R^{m\times n}$, $b\in\R^m$, and $c\in\R^n$, and read inequalities between vectors coordinatewise. The \term{primal} and the \term{dual} linear programs are
\begin{align*}
\text{(P)}\quad&\max\{c^\top x: Ax\le b,\ x\ge0\},\\
\text{(D)}\quad&\min\{b^\top y: A^\top y\ge c,\ y\ge0\}.
\end{align*}
A vector satisfying the constraints of a program is \term{feasible} for it, and a feasible vector attaining the maximum, or the minimum, is \term{optimal}.
\end{definition}

\begin{lemma}[Weak duality]\label{lem:lpduality:weak}
If $x$ is feasible for \textup{(P)} and $y$ is feasible for \textup{(D)}, then $c^\top x\le b^\top y$.
\end{lemma}

\begin{proof}
Use $A^\top y\ge c$ with $x\ge0$, and then $Ax\le b$ with $y\ge0$: $c^\top x\le (A^\top y)^\top x=y^\top(Ax)\le y^\top b$.
\end{proof}

Weak duality is every easy half at once.\index{duality!weak} For K\H{o}nig's theorem, let $M\in\{0,1\}^{V\times E}$ be the vertex--edge incidence matrix of the graph, with $M_{ve}=1$ when $v$ is an endpoint of $e$, and let $\lpone$ denote an all-ones vector. With $A=M$ and $b=c=\lpone$, program (P) asks for the largest $\sum_e x_e$ subject to weight at most $1$ at each vertex, a \term{fractional matching}, and (D) asks for the smallest $\sum_v y_v$ subject to $y_u+y_v\ge1$ on each edge, a \term{fractional vertex cover}. Incidence vectors of matchings and of covers are feasible, so the optimal values $\nu^*$ and $\tau^*$, which exist by \cref{thm:lpduality:duality} below, satisfy
\begin{equation}\label{eq:lpduality:chain}
\nu(G)\le\nu^*\le\tau^*\le\tau(G).
\end{equation}
The middle inequality is weak duality, and K\H{o}nig's theorem says that the whole chain collapses. The chain splits the theorem into two claims of different kinds. For the triangle the four numbers are $1$, $\frac32$, $\frac32$, and $2$: weight $\frac12$ on each edge is a fractional matching, weight $\frac12$ on each vertex is a fractional cover, and both have total $\frac32$, so both are optimal. The middle gap closes for every graph, by a theorem about all linear programs; the outer gaps are where bipartiteness matters.

\begin{theorem}[Duality theorem of linear programming]\label{thm:lpduality:duality}
If \textup{(P)} and \textup{(D)} both have feasible solutions, then both have optimal solutions, and their optimal values are equal.
\end{theorem}

The outer gaps are closed by a property of the matrix.\index{duality!strong}

\begin{definition}[Polyhedra, vertices, total unimodularity]\label{def:lpduality:tu}
Let $A\in\R^{m\times n}$ have rows $a_1^\top,\dots,a_m^\top$ and let $b\in\R^m$. The set $P=\{x\in\R^n:Ax\le b\}$ is a \term{polyhedron}, and a \term{polytope} if it is bounded and nonempty. The constraint $i$ is \term{tight} at $x\in P$ if $a_i^\top x=b_i$, and $x$ is a \term{vertex} of $P$ if the rows tight at $x$ span $\R^n$. A matrix is \term{totally unimodular} if each of its square submatrices has determinant $0$, $1$, or $-1$.
\end{definition}

\begin{theorem}[Integrality]\label{thm:lpduality:integrality}
If $A$ is totally unimodular and $b$ is integral, then every vertex of $\{x:Ax\le b\}$ is integral.
\end{theorem}

\Cref{thm:lpduality:duality,thm:lpduality:integrality} are proved in \cref{sec:lpduality:paid}.\sidenote{\textcite{HoffmanKruskal1956} proved a converse of \cref{thm:lpduality:integrality}: an integral matrix $A$ for which $\{x:Ax\le b,\ x\ge0\}$ has only integral vertices for every integral $b$ is totally unimodular.} Together they give a recipe: write the maximization as an integer program, drop integrality, take the dual, show that both programs have integral optima, and read the equality of optimal values as a statement about combinatorial objects. In one sentence, \emph{a min--max theorem is the duality theorem of linear programming applied to a pair of programs whose optima are integral.}\index{totally unimodular matrix}

The duality theorem predates its name. When George Dantzig described his new linear programming problem to von Neumann in October 1947, von Neumann answered with duality and conjectured its equivalence with his theory of games \parencite{Dantzig1982,vonStengel2024}.\sidenote{The minimax theorem follows easily from LP duality (\cref{sec:lpduality:zoo}). Dantzig's 1951 converse needs an extra assumption; \textcite{Adler2013} and \textcite{vonStengel2024} remove it.} The first published proof is due to \textcite{GaleKuhnTucker1951}. Its algebraic engine, a theorem of the alternative for linear inequalities, goes back to \textcite{Farkas1902}, and \textcite{HoffmanKruskal1956} identified total unimodularity as the condition under which integral right-hand sides force integral vertices.

In hindsight the move is natural, because a min--max theorem is a promise of short certificates, as in \cref{fig:lpduality:konig}, and linear programs always come with certificate pairs. The object that makes this visible is humble, a number between $0$ and $1$ on each edge. Nobody wants a half-matching, but fractional matchings form a convex polytope, and convex sets are where separation theorems live. What a reader of the zoo had to notice was that every easy half is the same inequality; it is then natural to ask whether every hard half is the same theorem.

\section{After: one duality theorem, paid for once}\label{sec:lpduality:paid}

This section proves \cref{thm:lpduality:duality,thm:lpduality:integrality}. The analytic content sits in two lemmas about closed convex sets; the rest is linear algebra.\index{separating hyperplane}

\begin{lemma}[Separation]\label{lem:lpduality:separation}
Let $K\subseteq\R^m$ be a nonempty closed convex set and let $b\in\R^m\setminus K$. Then some $y\in\R^m$ satisfies $y^\top b<y^\top z$ for every $z\in K$.
\end{lemma}

\begin{proof}
Fix $z_0\in K$. The continuous function $z\mapsto\norm{z-b}$ attains its minimum on the compact nonempty set $\{z\in K:\norm{z-b}\le\norm{z_0-b}\}$, at a point $z^*$ nearest to $b$ in all of $K$. Let $y=z^*-b$, which is nonzero because $b\notin K$. For $z\in K$ and $t\in(0,1]$, the point $z^*+t(z-z^*)$ lies in $K$ by convexity and is no nearer to $b$ than $z^*$. Expanding the squared distances gives
\[
2t\,(z-z^*)^\top y+t^2\norm{z-z^*}^2\ge0 .
\]
Divide by $t$ and let $t\to0$: then $(z-z^*)^\top y\ge0$, so $y^\top z\ge y^\top z^*=y^\top b+\norm{y}^2>y^\top b$.
\end{proof}

For $A\in\R^{m\times n}$ with columns $a_1,\dots,a_n$, the set $\lpcone(A)=\{Ax:x\ge0\}$ is a convex cone containing $0$. Separation needs it to be closed.

\begin{lemma}[Finitely generated cones are closed]\label{lem:lpduality:cone}
For every matrix $A$, the cone $\lpcone(A)$ is closed.
\end{lemma}

\begin{proof}
The plan is to write the cone as a finite union of cones generated by linearly independent columns, each of which is closed.

First, every $z\in\lpcone(A)$ is a nonnegative combination of linearly independent columns. Take a representation $z=Ax$ with $x\ge0$ and as few nonzero entries as possible, and let $J$ be the set of indices $j$ with $x_j>0$. If the columns $a_j$ with $j\in J$ were dependent, then some nonzero $\mu$ supported on $J$ would have $A\mu=0$ and, after a change of sign if necessary, a positive entry $\mu_j$. With $\theta=\min\{x_j/\mu_j:\mu_j>0\}$, the vector $x-\theta\mu$ would be nonnegative, represent $z$, and have fewer nonzero entries, a contradiction.

Second, $\{Bw:w\ge0\}$ is closed if $B$ consists of linearly independent columns of $A$. Indeed $B^\top B$ is invertible, and if $Bw_k\to z$ with $w_k\ge0$, then $w_k=(B^\top B)^{-1}B^\top(Bw_k)$ converges to some $w$, which is nonnegative and satisfies $Bw=z$. So $\lpcone(A)$ is the union of finitely many closed sets, one for each set of linearly independent columns, and is closed.
\end{proof}

\begin{theorem}[\textcite{Farkas1902}]\label{thm:lpduality:farkas}
Let $A\in\R^{m\times n}$ and $b\in\R^m$.
\begin{enumerate}
\item Exactly one of the following holds: some $x\ge0$ satisfies $Ax=b$; or some $y$ satisfies $y^\top A\ge0$ and $y^\top b<0$.
\item Exactly one of the following holds: some $x\ge0$ satisfies $Ax\le b$; or some $y\ge0$ satisfies $y^\top A\ge0$ and $y^\top b<0$.
\end{enumerate}
\end{theorem}

\begin{proof}
(i) Both cannot hold, because $Ax=b$ with $x\ge0$ and $y^\top A\ge0$ give $y^\top b=(y^\top A)x\ge0$. Suppose the first fails, so that $b\notin\lpcone(A)$. The cone is nonempty, convex, and closed (\cref{lem:lpduality:cone}), so \cref{lem:lpduality:separation} provides $y$ with $y^\top b<y^\top z$ for every $z$ in the cone. Taking $z=0$ gives $y^\top b<0$. Taking $z=\lambda a_j$ with $\lambda>0$ gives $y^\top a_j>y^\top b/\lambda$, and letting $\lambda\to\infty$ gives $y^\top a_j\ge0$ for every column.

(ii) Some $x\ge0$ satisfies $Ax\le b$ exactly when some $(x,s)\ge0$ satisfies $Ax+s=b$. Apply (i) to the matrix $[A\ \ I]$, and note that $y^\top[A\ \ I]\ge0$ says that $y^\top A\ge0$ and $y\ge0$.
\end{proof}\index{Farkas's lemma}

\begin{figure}[t]
\centering
\begin{tikzpicture}[scale=0.82,>=Stealth]
  \fill[seamist] (0,0)--(3.2,0.8)--(3.2,3.0)--(1.0,3.0)--cycle;
  \draw[sandline,line width=0.8pt] (0.66,-2.2)--(-0.9,3.0);
  \draw[->,sea,line width=1pt] (0,0)--(2,0.5) node[below right=-1pt,font=\small] {$a_1$};
  \draw[->,sea,line width=1pt] (0,0)--(0.6,1.8) node[left=1pt,font=\small] {$a_2$};
  \draw[->,seadeep,line width=1pt] (0,0)--(1.0,0.3) node[below=2pt,font=\small] {$y$};
  \fill[seadeep] (-1.2,1.4) circle (1.6pt) node[left=2pt,font=\small] {$b$};
  \fill[ink] (0,0) circle (1pt);
  \node[font=\small,text=seadeep] at (2.45,2.35) {$\lpcone(A)$};
  \node[font=\footnotesize,text=muted] at (-1.9,-0.6) {$y^\top z<0$};
  \node[font=\footnotesize,text=muted] at (1.9,-0.9) {$y^\top z\ge 0$};
\end{tikzpicture}
\caption{Farkas's lemma in the plane. A point $b$ outside the cone generated by $a_1$ and $a_2$ is separated from it by a line through the origin, whose normal $y$ has $y^\top a_j\ge0$ for both columns and $y^\top b<0$. The vector $y$ certifies that $Ax=b$ has no solution $x\ge0$.}\label{fig:lpduality:farkas}
\end{figure}

\begin{proof}[Proof of \cref{thm:lpduality:duality}]
Let $\bar x$ and $\bar y$ be feasible for (P) and (D). By weak duality it suffices to find feasible $x$ and $y$ with $c^\top x\ge b^\top y$, that is, a nonnegative solution $(x,y)$ of the system
\[
Ax\le b,\qquad -A^\top y\le -c,\qquad -c^\top x+b^\top y\le 0 .
\]
Suppose there is none. \Cref{thm:lpduality:farkas}(ii), applied to this system, gives nonnegative multipliers $u\in\R^m$, $v\in\R^n$, and $\lambda\in\R$ for its three blocks with
\[
A^\top u\ge\lambda c,\qquad Av\le\lambda b,\qquad b^\top u<c^\top v .
\]
If $\lambda>0$, then $v/\lambda$ is feasible for (P) and $u/\lambda$ is feasible for (D), and $b^\top u<c^\top v$ contradicts weak duality. If $\lambda=0$, then $A^\top u\ge0$ and $Av\le0$. Use $A\bar x\le b$ with $u\ge0$, then $\bar x\ge0$ with $A^\top u\ge0$, then $\bar y\ge0$ with $Av\le0$, and finally $A^\top\bar y\ge c$ with $v\ge0$:
\[
b^\top u\ge\bar x^\top A^\top u\ge0\ge\bar y^\top Av\ge c^\top v ,
\]
which again contradicts $b^\top u<c^\top v$. So the system has a nonnegative solution, and weak duality forces $c^\top x=b^\top y$: both $x$ and $y$ are optimal, with equal values.
\end{proof}

An equality constraint $a^\top x=\beta$ is the pair $a^\top x\le\beta$ and $-a^\top x\le-\beta$, whose two nonnegative multipliers enter (D) only through their difference. So equalities in one program correspond to variables of either sign in the other, and \cref{thm:lpduality:duality} holds verbatim for such programs.

Integrality needs two more lemmas: one about where linear functions attain their maxima, and one about which matrices are totally unimodular.\sidenote{For compact convex sets the analogue of \cref{lem:lpduality:vertices} is Minkowski's theorem in $\R^n$, with extreme points for vertices, and in locally convex spaces the Krein--Milman theorem, which takes closed convex hulls.}\index{vertex of a polyhedron}

\begin{lemma}[Vertices]\label{lem:lpduality:vertices}
Every polytope $P=\{x:Ax\le b\}$ has finitely many vertices, and every point of $P$ is a convex combination of them. Consequently every linear function attains its maximum over $P$ at a vertex.
\end{lemma}

\begin{proof}
A vertex is the unique solution of a nonsingular square system formed by $n$ of its tight constraints, and there are finitely many such systems. For $x\in P$ let $r(x)$ be the dimension of the span of the rows tight at $x$. We show by downward induction on $r(x)$ that $x$ is a convex combination of vertices. If $r(x)=n$, then $x$ is a vertex. Otherwise some nonzero $d$ is orthogonal to every row tight at $x$. The other constraints have slack and $P$ is bounded, so there are largest numbers $t^+>0$ and $t^->0$ with $y=x+t^+d$ and $z=x-t^-d$ in $P$. The rows tight at $x$ stay tight at $y$, and at least one new row $a_i$ with $a_i^\top d\ne0$ becomes tight; it is not in the span of the rows tight at $x$, which are orthogonal to $d$. So $r(y)>r(x)$, and likewise $r(z)>r(x)$. By induction $y$ and $z$ are convex combinations of vertices, and so is $x=(t^-y+t^+z)/(t^++t^-)$. Finally, if $x=\sum_k\lambda_kv_k$ is a convex combination of vertices, then $c^\top x\le\max_kc^\top v_k$, so the largest value of $c^\top$ on the finitely many vertices is its maximum over $P$.
\end{proof}

\begin{proof}[Proof of \cref{thm:lpduality:integrality}]
Let $x$ be a vertex. Choose $n$ linearly independent rows tight at $x$; they form a nonsingular $n\times n$ submatrix $B$ of $A$ with $Bx=b_B$, where $b_B$ is the corresponding part of $b$. By Cramer's rule, $x_j$ is the determinant of an integral matrix ($B$ with its $j$th column replaced by $b_B$) divided by $\det B$, which is $\pm1$ by total unimodularity. So $x$ is integral.
\end{proof}

\begin{lemma}[Building totally unimodular matrices]\label{lem:lpduality:incidence}
\leavevmode
\begin{enumerate}
\item If $A$ is totally unimodular, then so are $A^\top$ and every matrix obtained from $A$ by repeatedly negating a row, repeating a row, or appending a row of an identity matrix.
\item A matrix with entries in $\{0,1,-1\}$ in which each column has at most one entry $1$ and at most one entry $-1$ is totally unimodular. In particular, the incidence matrix of a loopless directed graph (whose column for an arc from $u$ to $v$ has $-1$ in row $u$ and $1$ in row $v$), every submatrix of it, and the vertex--edge incidence matrix of a bipartite graph are totally unimodular.
\end{enumerate}
\end{lemma}

\begin{proof}
(i) Transposing and negating rows change determinants at most in sign. A square submatrix that uses a row twice, or a row and its negation, has determinant $0$. A square submatrix that uses a row of the identity either has a zero row or can be expanded along that row, which has a single entry $\pm1$; induction on the size reduces to a square submatrix of $A$.

(ii) We induct on the size of a square submatrix $B$, which inherits the hypothesis on columns. Size $1$ is clear. A zero column gives $\det B=0$. If some column has exactly one nonzero entry, then expanding along it reduces the size by one. Otherwise each column of $B$ has one entry $1$ and one entry $-1$, so the rows of $B$ sum to zero and $\det B=0$. For a bipartite graph with sides $L$ and $R$, negate the rows of $R$: each column then has $1$ in a row of $L$ and $-1$ in a row of $R$, and (i) undoes the negation.
\end{proof}

\section{After: the zoo, re-proved}\label{sec:lpduality:zoo}

With the price paid, each member of the zoo costs a formulation and a check.

\begin{proof}[Second proof of \cref{thm:lpduality:konig}]
Let $M$ be the vertex--edge incidence matrix of the bipartite graph; it is totally unimodular (\cref{lem:lpduality:incidence}). Programs (P) and (D) of \eqref{eq:lpduality:chain} are feasible ($x=0$ and $y=\lpone$), so they have optimal solutions with a common value $\nu^*=\tau^*$ (\cref{thm:lpduality:duality}). The feasible set of (P) is a polytope, since $0\le x_e\le1$, with totally unimodular constraint matrix $\left[\begin{smallmatrix}M\\-I\end{smallmatrix}\right]$; by \cref{lem:lpduality:vertices} and \cref{thm:lpduality:integrality} it has an integral optimal vertex, the incidence vector of a matching, so $\nu=\nu^*$. For (D), replacing an optimal $y$ by $\min(y,\lpone)$ keeps every constraint $y_u+y_v\ge1$ and does not raise the cost, so $\tau^*$ is also the minimum over the polytope $\{y:M^\top y\ge\lpone,\ 0\le y\le\lpone\}$, whose matrix is again totally unimodular. An integral optimal vertex is the incidence vector of a vertex cover, so $\tau=\tau^*$.
\end{proof}

\begin{proof}[Second proof of \cref{thm:lpduality:bvn}]
The doubly stochastic $n\times n$ matrices form the polytope $\{X\ge0: X\lpone=\lpone,\ X^\top\lpone=\lpone\}$. Its equality constraints have as matrix the incidence matrix of the complete bipartite graph on rows and columns, so all its vertices are integral (\cref{lem:lpduality:incidence}, \cref{thm:lpduality:integrality}). An integral doubly stochastic matrix is a permutation matrix, and every point of the polytope is a convex combination of vertices (\cref{lem:lpduality:vertices}).
\end{proof}

The induction in the first proof of \cref{thm:lpduality:bvn} did not disappear; it moved into \cref{lem:lpduality:vertices}, done once for all polytopes.

Next, flows. Let $D=(V,A)$ be a loopless directed graph with distinct vertices $s$ and $t$ and capacities $c\in\R^A$, $c\ge0$. A \term{flow} is a vector $f\in\R^A$ with $0\le f\le c$ such that at every vertex other than $s$ and $t$ the flow on entering arcs equals the flow on leaving arcs. Its \term{value} $\lpval(f)$ is the flow on arcs leaving $s$ minus the flow on arcs entering $s$. A \term{cut} is a set $S\subseteq V$ with $s\in S$ and $t\notin S$; its capacity $c(S)$ is the total capacity of the arcs leaving $S$.\index{max-flow min-cut theorem}

\begin{theorem}[Max-flow min-cut]\label{thm:lpduality:maxflow}
The largest flow value equals the smallest cut capacity. If $c$ is integral, then some maximum flow is integral.
\end{theorem}

\begin{proof}
The plan follows the second proof of K\H{o}nig's theorem, with a rounding step in place of truncation on the dual side.

\emph{The easy half.} For a flow $f$ and a cut $S$, add the conservation equations of the vertices of $S\setminus\{s\}$ to the definition of $\lpval(f)$. The arcs with both ends in $S$ cancel, and what remains is the flow on arcs leaving $S$ minus the flow on arcs entering $S$, which is at most $c(S)$.

\emph{The maximum.} Let $N$ be the incidence matrix of $D$ with the rows of $s$ and $t$ deleted. The flows form the polytope $\{f:Nf=0,\ 0\le f\le c\}$, which contains $f=0$ and has a totally unimodular constraint matrix (\cref{lem:lpduality:incidence}). By \cref{lem:lpduality:vertices} the value attains its maximum $\varphi$ at a vertex, which is integral when $c$ is (\cref{thm:lpduality:integrality}).

\emph{The cut.} The value is $w^\top f$, where $w_a$ is $1$ on arcs leaving $s$, $-1$ on arcs entering $s$, and $0$ elsewhere. Give the conservation constraint of each vertex $v\ne s,t$ a multiplier $\pi_v$ of either sign and the capacity constraint of each arc $a$ a multiplier $z_a\ge0$. The dual constraint of an arc $a$ from $u$ to $v$ is $z_a\ge w_a+\pi_u-\pi_v$, where $\pi_s=\pi_t=0$; setting $\pi_s=1$ instead turns it into $z_a\ge\pi_u-\pi_v$, so the dual program is
\[
\min\Bigl\{\textstyle\sum_a c_az_a:\ z\ge0,\ z_a\ge\pi_u-\pi_v\text{ for every arc }a\text{ from }u\text{ to }v\Bigr\}.
\]
It is feasible, so \cref{thm:lpduality:duality} gives a feasible $(\pi,z)$ of cost $\varphi$. Choose $\theta$ uniformly at random from $(0,1]$ and let $S_\theta=\{v:\pi_v\ge\theta\}$, a cut because $\pi_s=1$ and $\pi_t=0$. An arc from $u$ to $v$ leaves $S_\theta$ only when $\pi_v<\theta\le\pi_u$, an event of probability at most $\max(\pi_u-\pi_v,0)\le z_a$. So the expected capacity of $S_\theta$ is at most $\varphi$, some cut has capacity at most $\varphi$, and the easy half finishes the proof.
\end{proof}

With unit capacities an integral flow of value $k$ splits into $k$ arc-disjoint paths from $s$ to $t$ (a routine decomposition that we omit), and \cref{thm:lpduality:maxflow} becomes the arc version of Menger's theorem; splitting each vertex into an arc gives the vertex version. Dilworth's theorem needs no new program at all: \textcite{Fulkerson1956} reduced it to K\H{o}nig's (\cref{ex:lpduality:dilworth}).\index{Dilworth's theorem}

Finally, the member from game theory.

\begin{proof}[Proof of \cref{thm:lpduality:minimax} from LP duality]
For a fixed $p$, the minimum of $p^\top Aq$ over $q\in\Delta_n$ is the smallest entry of $p^\top A$, a continuous function of $p$, so the maximum on the left exists; likewise on the right.

\emph{The easy half.} For all $p$ and $q$, $\min_{q'}p^\top Aq'\le p^\top Aq\le\max_{p'}p'^\top Aq$. Maximize over $p$ on the left and minimize over $q$ on the right.

\emph{The hard half.} Adding a constant to every entry of $A$ adds it to $p^\top Aq$ for all mixed strategies, so we may assume that every entry of $A$ is positive. The programs
\[
\max\{\lpone^\top u: Au\le\lpone,\ u\ge0\}\qquad\text{and}\qquad\min\{\lpone^\top w: A^\top w\ge\lpone,\ w\ge0\}
\]
form a pair as in \cref{def:lpduality:pair}, and both are feasible: take $u=0$, and $w$ equal to the first unit vector divided by the smallest entry of the first row of $A$. By \cref{thm:lpduality:duality} they have optimal solutions $u^*$ and $w^*$ with $\lpone^\top u^*=\lpone^\top w^*=\theta$, and $\theta>0$ because $A^\top w^*\ge\lpone$ forces $w^*\ne0$. The vectors $q^*=u^*/\theta$ and $p^*=w^*/\theta$ are mixed strategies. From $Aq^*\le\lpone/\theta$, every $p\in\Delta_m$ has $p^\top Aq^*\le1/\theta$; from $A^\top p^*\ge\lpone/\theta$, every $q\in\Delta_n$ has $p^{*\top}Aq\ge1/\theta$. So the left side is at least $1/\theta$, and the right side at most $1/\theta$.
\end{proof}

Weighted versions cost nothing extra. With integral weights $w\ge0$ as the objective of (P), and the dual truncated at $\max_ew_e$ instead of $1$, the second proof of K\H{o}nig's theorem gives Egerv\'ary's theorem.

\section{After: duality for algorithms, learning, and hardness}\label{sec:lpduality:reach}

Read von Neumann's theorem as a statement about responses: \emph{if every mixed strategy of one player has a good pure response, then some single mixed strategy of the other player is good against everything.} Three statements of theoretical computer science have this form without mentioning a game. In the last two, the minimax theorem produces a mixed strategy, and sampling turns it into a small object.

\begin{lemma}[Majority of samples]\label{lem:lpduality:sampling}
Let $(x_1,y_1),\dots,(x_N,y_N)$ be points with labels in a two-element set, let $\lpQ$ be a finitely supported probability distribution on functions taking values in that set, and let $\eta>0$. Suppose that $\Pr_{h\sim\lpQ}[h(x_i)=y_i]\ge\frac12+\eta$ for every $i$. Let $h_1,\dots,h_t$ be independent samples from $\lpQ$, and let $g$ be their majority vote, with ties broken arbitrarily. Then $\Pr[g(x_i)\ne y_i]\le e^{-2t\eta^2}$ for every $i$. Consequently, some choice of $h_1,\dots,h_t$ in the support of $\lpQ$ makes $g$ err on at most $Ne^{-2t\eta^2}$ of the points, and if $t>\ln N/(2\eta^2)$, then some choice makes $g$ err on none.
\end{lemma}

\begin{proof}
Fix $i$. The indicators of the events $h_j(x_i)=y_i$ are independent with mean at least $\frac12+\eta$, and $g(x_i)=y_i$ whenever more than half of them equal~$1$. Hoeffding's inequality says that an average of $t$ independent random variables with values in $[0,1]$ falls at least $\eta$ below its mean with probability at most $e^{-2t\eta^2}$, so their average is at most $\frac12$ with at most this probability. By linearity of expectation, the expected number of points on which $g$ errs is at most $Ne^{-2t\eta^2}$, so some outcome of the samples achieves at most this number; when $t>\ln N/(2\eta^2)$, this number is less than $1$.
\end{proof}

\subsection{Yao's principle}

Let $X$ be a finite set of inputs and $\cA$ a set of deterministic algorithms for some task, with a cost $\lpcost(A,x)\ge0$ of running $A$ on $x$. A randomized algorithm is a probability distribution $R$ on $\cA$: flip all the coins in advance, and a deterministic algorithm remains.\index{Yao's principle}

\begin{theorem}[\textcite{Yao1977}]\label{thm:lpduality:yao}
\leavevmode
\begin{enumerate}
\item For every randomized algorithm $R$ and every probability distribution $D$ on inputs,
\[
\max_{x\in X}\E_{A\sim R}\lpcost(A,x)\ \ge\ \inf_{A\in\cA}\E_{x\sim D}\lpcost(A,x).
\]
\item If $\cA$ is finite, then the smallest worst-case expected cost of a randomized algorithm equals the largest, over distributions $D$ on inputs, of the smallest expected cost against $D$ of a deterministic algorithm.
\end{enumerate}
\end{theorem}

\begin{proof}
(i) is weak duality: the maximum over $x$ is at least the average over $x\sim D$; the two expectations may be exchanged because $X$ is finite; and an average over $A\sim R$ is at least the infimum over $A$. (ii) is \cref{thm:lpduality:minimax} for the matrix with entry $\lpcost(A,x)$ in row $x$ and column $A$; against a fixed mixed strategy some pure strategy is a best response, so the inner optimizations may run over single inputs and single algorithms.
\end{proof}

Part (i) is all that a lower bound needs, and it holds for infinite sets of algorithms;\sidenote{In statistics the same inequality says that the Bayes risk of every prior bounds the minimax risk from below; \cref{ch:lecam} develops lower bounds of this kind.} part (ii) says that some input distribution always proves the optimal lower bound. For an example, hide an item in one of $n$ cells; an algorithm opens cells until it finds the item and pays the number of cells opened. A deterministic algorithm opens the cells in a fixed order, since every cell it opens before success is empty, so it pays at least $n$ on some input, while a uniformly random order pays $(n+1)/2$ in expectation on every input. Against uniformly placed items every deterministic algorithm pays at least $(n+1)/2$ on average, so by part (i) no randomized algorithm does better: the random order and the uniform placement certify each other, as the matching and the cover of \cref{fig:lpduality:konig} do.

\subsection{Weak learning is strong learning on a sample}

A weak learner finds, on every reweighting of the data, a hypothesis slightly better than a coin flip. Kearns and Valiant asked whether weak learnability implies learnability to arbitrary accuracy; \textcite{Schapire1990} answered yes with an intricate recursive construction (\cref{ch:regret}), and \textcite{FreundSchapire1996} observed that on a fixed sample the existence half of the answer is von Neumann's theorem.\index{boosting}\index{weak learning}

Let $(x_1,y_1),\dots,(x_m,y_m)$ be examples with labels $y_i\in\{-1,1\}$, and let $\cH$ be a nonempty set of functions with values in $\{-1,1\}$. For a distribution $D$ on $\{1,\dots,m\}$, the weighted error of $h$ is $\lperr_D(h)=\sum_iD_i\ind[h(x_i)\ne y_i]$.

\begin{theorem}[Weak learning on a sample]\label{thm:lpduality:boosting}
Let $\gamma\in(0,\frac12]$. The following are equivalent.
\begin{enumerate}
\item For every distribution $D$ on the examples, some $h\in\cH$ has $\lperr_D(h)\le\frac12-\gamma$.
\item Some finitely supported distribution $\lpQ$ on $\cH$ satisfies $\sum_h\lpQ(h)\,y_ih(x_i)\ge2\gamma$ for every $i$.
\end{enumerate}
If they hold, then for every integer $t>\ln m/(2\gamma^2)$ some $h_1,\dots,h_t\in\cH$ have a majority vote that is correct on all $m$ examples.
\end{theorem}

\begin{proof}
Because $y_ih(x_i)\in\{-1,1\}$, the indicator of an error is $(1-y_ih(x_i))/2$, so $\lperr_D(h)\le\frac12-\gamma$ means $\sum_iD_iy_ih(x_i)\ge2\gamma$. Hypotheses that agree on $x_1,\dots,x_m$ are interchangeable in (i) and (ii), so we may keep one hypothesis for each of the at most $2^m$ patterns of values and assume $\cH$ finite. Let $W$ be the $m\times|\cH|$ matrix with entries $W_{ih}=y_ih(x_i)$. Statement (i) says that $\min_D\max_h(D^\top W)_h\ge2\gamma$, and (ii) says that $\max_{\lpQ}\min_i(W\lpQ)_i\ge2\gamma$. The two numbers are equal by \cref{thm:lpduality:minimax}, applied to $W^\top$ with the hypothesis player maximizing. For the last claim, (ii) gives $\Pr_{h\sim\lpQ}[h(x_i)=y_i]=\frac12+\frac12\sum_h\lpQ(h)y_ih(x_i)\ge\frac12+\gamma$ for every $i$, and \cref{lem:lpduality:sampling} with $\eta=\gamma$ and $N=m$ applies.
\end{proof}

The direction from (ii) to (i) is weak duality.\sidenote{AdaBoost runs Hedge, with a varying rate, on the examples against the weak learner; its weighted vote fits the sample after more than $\ln m/(2\gamma^2)$ rounds. With a fixed rate, Hedge finds for every $\eps>0$ a mixture as in (ii) with $2\gamma-\eps$ in place of $2\gamma$. See \cref{ch:regret}.} The other direction says that if weak learning succeeds on every reweighting of the sample, then a majority of $O(\log m/\gamma^2)$ weak hypotheses fits the whole sample; predicting well on new data needs a further argument, through margins or compression.

\subsection{The hardcore lemma}

A Boolean function $f:\{0,1\}^n\to\{0,1\}$ may be mildly hard for small circuits, each of them disagreeing with $f$ on at least a $\delta$ fraction of inputs. \textcite{Impagliazzo1995} proved that such hardness concentrates on a \term{hardcore} distribution of density $\delta$, where no small circuit does much better than a coin; the minimax proof below is Nisan's, reported in Impagliazzo's paper. A distribution $H$ on $\{0,1\}^n$ has \term{density} $\delta$ if $H(x)\le1/(\delta2^n)$ for every $x$.\index{hardcore lemma}

\begin{theorem}[Hardcore lemma]\label{thm:lpduality:hardcore}
Let $f:\{0,1\}^n\to\{0,1\}$, let $\cC$ be a finite nonempty set of functions $\{0,1\}^n\to\{0,1\}$, let $k\in\{1,\dots,2^n\}$ and $\delta=k/2^n$, let $\eps\in(0,1)$, and let $t\ge(2/\eps^2)\ln(2/(\eps\delta))$ be an integer. Suppose that every majority vote of $t$ members of $\cC$, not necessarily distinct, disagrees with $f$ on at least $k$ inputs. Then some distribution $H$ of density $\delta$ satisfies $\Pr_{x\sim H}[C(x)=f(x)]\le\frac12+\eps$ for every $C\in\cC$.
\end{theorem}

\begin{proof}
The plan is to suppose the conclusion false, obtain from the minimax theorem a mixture of members of $\cC$ that predicts $f$ well on every set of $k$ inputs, and sample from it a majority vote that contradicts the hypothesis. Let a predictor choose $C\in\cC$ and an adversary a set $S$ of $k$ inputs, the predictor receiving the fraction of $S$ on which $C$ agrees with $f$. A mixed strategy of the adversary induces a mixture of uniform distributions on $k$-sets, which has density $\delta$. If the conclusion fails, then each such mixture is beaten by some $C$ with payoff above $\frac12+\eps$, so the value of the game, attained at some mixed strategy of the adversary, exceeds $\frac12+\eps$, and by \cref{thm:lpduality:minimax} some mixed strategy $\lpQ$ of the predictor gets more than $\frac12+\eps$ against every $k$-set: the average over every set of $k$ inputs of $a(x)=\Pr_{C\sim\lpQ}[C(x)=f(x)]$ exceeds $\frac12+\eps$.

Let $B=\{x:a(x)\le\frac12+\frac\eps2\}$. If $|B|\ge(1-\eps)k$, then a $k$-set containing $\min(|B|,k)$ points of $B$ has an average of $a$ at most $1-(1-\eps)(\frac12-\frac\eps2)<\frac12+\eps$, since points outside $B$ contribute at most $1$; so $|B|<(1-\eps)k$. Every $x\notin B$ has $a(x)>\frac12+\frac\eps2$, so \cref{lem:lpduality:sampling} with $\eta=\eps/2$ gives $t$ members of $\cC$ whose majority vote errs on at most $2^ne^{-t\eps^2/2}\le\eps k/2$ points outside $B$, by the choice of $t$. The vote disagrees with $f$ on fewer than $(1-\eps)k+\eps k/2<k$ inputs, contradicting the hypothesis.
\end{proof}

For circuits, let $\cC$ be the circuits of size at most $s'$; a majority of $t$ of them has size $ts'+O(t)$, so the loss in circuit size is a factor of order $\eps^{-2}\log(1/(\eps\delta))$. Impagliazzo's own argument was constructive, an iteration resembling boosting, and the lemma is one standard route to Yao's XOR lemma (\cref{ch:pseudorandom}).

\Cref{thm:lpduality:boosting,thm:lpduality:hardcore} have one proof: minimax, then sampling. In both, the minimax theorem asserts that a good mixture exists and gives no way to find it. \Cref{ch:regret} proves the minimax theorem by dynamics instead, and these existence statements become algorithms; for weak learning, the algorithm is boosting.

\section{Costs and limits}\label{sec:lpduality:costs}

\paragraph{Where did the difficulty go?} Into two places. The analytic content of every min--max theorem in this chapter is the separation step: the nearest point of \cref{lem:lpduality:separation}, which needs compactness, and the closedness of cones in \cref{lem:lpduality:cone}. The combinatorial content went into total unimodularity, a determinant condition checked once for each family of matrices (\cref{lem:lpduality:incidence}). After both payments a new min--max theorem, and its weighted version, cost a formulation and a check; that is why the bargain is real. Something was given up, though: the augmenting-path proof of K\H{o}nig's theorem is an algorithm, while the duality proof only asserts that a certificate exists. Finding one is a separate problem, solved by the simplex method in practice and in polynomial time by the ellipsoid and interior-point methods \parencite{Schrijver1986}; in \cref{sec:lpduality:reach} the minimax theorem names no procedure at all.

\paragraph{What did not dissolve.} Matching in general graphs. For the triangle, the chain \eqref{eq:lpduality:chain} reads $1<\frac32=\frac32<2$: duality closes the middle gap, but the incidence matrix of an odd cycle is not totally unimodular and both outer gaps stay open. \textcite{Edmonds1965} found the extra inequalities, one for each odd set of vertices, that make the matching polytope integral, and the resulting theory goes beyond total unimodularity. Hard problems presumably have no theorem of K\H{o}nig's type: a min--max equality with certificates that are easy to check puts a decision problem in NP and in coNP at once, so an NP-hard problem cannot have one unless NP equals coNP. And the minimax theorem needs finiteness or compactness. If each player names a positive integer and the larger number wins a unit, then every mixed strategy of either player can be answered by a number that beats it with probability as close to $1$ as we like; the best guarantees are $-1$ for one player and $1$ for the other, and no value exists.

\paragraph{Other seas.} Lagrangian and conic duality extend the move to convex programs, where strong duality needs a constraint qualification, and total dual integrality \parencite{Schrijver1986} yields min--max theorems where total unimodularity fails, among them matroid intersection (\cref{ch:matroid}).

\section{The move, distilled}\label{sec:lpduality:distilled}

\begin{move}{Duality}
\movepart{In one sentence} When a maximum keeps meeting a minimum, write both as linear programs: the easy half is weak duality, the hard half is strong duality, proved once, and integrality is a property of the constraint matrix.
\movepart{Triggers} A min--max statement whose easy half is a one-line count; a statement of the form ``for every distribution there is a good response''; fractional relaxations that seem to have the same optimum as the integral problem; a pair of certificates, each checking the other.
\movepart{Recipe} Write the maximization as an integer program; drop integrality; take the dual; prove that both programs have integral optima, by total unimodularity or by rounding; read off the combinatorial statement. For game-like statements, identify the bilinear payoff, pass to mixed strategies, apply the minimax theorem, and sparsify the resulting mixture by sampling.
\movepart{Anti-pattern} Dualizing a relaxation with an integrality gap and calling the result a min--max theorem; applying the minimax theorem to a payoff that is not linear in the mixed strategies, or to infinite strategy sets without compactness.
\end{move}

\dissolutiontest{K\H{o}nig's theorem in a paragraph instead of a page; Birkhoff's theorem in a few lines instead of an induction resting on K\H{o}nig's theorem.}{Weights and real capacities come free; bipartiteness is replaced by a determinant condition that says exactly what is needed.}{K\H{o}nig, Egerv\'ary, Menger, max-flow min-cut, Dilworth, Birkhoff--von Neumann, and von Neumann's minimax theorem become one duality theorem plus integrality.}{Yao's principle, weak learning implies strong learning on a sample, and the hardcore lemma.}

This move opens a chain. \Cref{ch:regret} keeps the minimax theorem and replaces its proof by dynamics, so that the existence statements of \cref{sec:lpduality:reach} become algorithms. \Cref{ch:matroid} meets duality in the matroid polytope, where greed constructs the dual solution, \cref{ch:transport} in Kantorovich's dual problem, \cref{ch:lecam} in lower bounds through priors, and \cref{ch:sos} in a duality between proofs and pseudo-distributions.

\begin{exercises}
\item \consolidation Prove Hall's theorem from K\H{o}nig's theorem: a bipartite graph with sides $L$ and $R$ has a matching that touches every vertex of $L$ if and only if every set $S\subseteq L$ has at least $|S|$ neighbors.\label{ex:lpduality:hall}

\item \consolidation (a) Prove that $Ax\le b$ has no solution $x\in\R^n$ if and only if some $y\ge0$ satisfies $y^\top A=0$ and $y^\top b<0$. (b) Prove Gordan's theorem: exactly one of the following holds: $Ax=0$ for some nonzero $x\ge0$; or every entry of $y^\top A$ is positive for some $y$.\label{ex:lpduality:gordan}

\item \consolidation\hard Show that if (P) is feasible and its objective is bounded above on the feasible set, then (D) is feasible. Conclude that if one program of a pair has an optimal solution, then so does the other, with the same value. \hint{If (D) is infeasible, apply \cref{thm:lpduality:farkas}(ii) to the system $-A^\top y\le-c$, $y\ge0$, and move along the direction it provides.}\label{ex:lpduality:bounded}

\item \consolidation (a) Show that the vertex--edge incidence matrix of a cycle of odd length has a determinant of absolute value $2$. (b) Conclude that the vertex--edge incidence matrix of a loopless graph is totally unimodular if and only if the graph is bipartite.\label{ex:lpduality:odd}

\item \consolidation (Dilworth from K\H{o}nig.) Let $(X,<)$ be a finite partial order, and form the bipartite graph with sides $\{x':x\in X\}$ and $\{x'':x\in X\}$ and an edge $x'y''$ whenever $x<y$. (a) Show that a matching $\cM$ yields a partition of $X$ into $|X|-|\cM|$ chains: call $y$ the successor of $x$ when $x'y''\in\cM$. (b) Show that for a vertex cover $C$, the elements $x$ with neither $x'$ nor $x''$ in $C$ form an antichain of size at least $|X|-|C|$. (c) Deduce Dilworth's theorem.\label{ex:lpduality:dilworth}

\item \consolidation (Complementary slackness.) Let $x$ and $y$ be feasible for (P) and (D). Show that both are optimal if and only if $x_j>0$ implies $(A^\top y)_j=c_j$ and $y_i>0$ implies $(Ax)_i=b_i$. Interpret the conditions for a maximum matching and a minimum cover of a bipartite graph.\label{ex:lpduality:cs}

\item \transfer (Support vector machines.) Let points $x_1,\dots,x_m$ of $\R^d$ have labels $y_i\in\{-1,1\}$, and let $w$ satisfy $y_iw^\top x_i\ge1$ for every $i$. Show that every $\alpha\in\R^m$ with $\alpha\ge0$ satisfies $\frac12\norm{w}^2\ge\sum_i\alpha_i-\frac12\norm{\sum_i\alpha_iy_ix_i}^2$, with equality if and only if $w=\sum_i\alpha_iy_ix_i$ and $\alpha_i=0$ whenever $y_iw^\top x_i>1$. Points with $\alpha_i>0$ are support vectors.\label{ex:lpduality:svm}

\item \transfer (Kantorovich duality.) Let $\mu\in\Delta_m$, $\nu\in\Delta_n$, and $c\in\R^{m\times n}$. Prove that the smallest value of $\sum_{ij}c_{ij}\pi_{ij}$ over $\pi\ge0$ with $\pi\lpone=\mu$ and $\pi^\top\lpone=\nu$ equals the largest value of $\sum_i\mu_i\varphi_i+\sum_j\nu_j\psi_j$ over $\varphi,\psi$ with $\varphi_i+\psi_j\le c_{ij}$ for all $i,j$. For $m=n$ and uniform $\mu=\nu$, show that some optimal $\pi$ is a permutation matrix divided by $n$ (see \cref{ch:transport}).\label{ex:lpduality:kantorovich}

\item \transfer (Yao's principle for sorting.) Show that every randomized comparison sort of $n$ distinct keys that always sorts correctly makes, on some input order, at least $\log_2(n!)$ comparisons in expectation. \hint{Use the uniform distribution on orders and the inequality $\sum_i2^{-d_i}\le1$ for the depths $d_i$ of distinct leaves of a binary tree.}\label{ex:lpduality:sorting}

\item \research\harder (Regularity as duality.) Let $\cX$ be a finite set with the uniform distribution, $\cF$ a family of functions $\cX\to[-1,1]$, $g:\cX\to[0,1]$, and $\eps>0$. (a) Prove that some $k\le1/\eps^2$ functions $f_1,\dots,f_k\in\cF$ and some $h$ in their linear span satisfy $\abs{\E[f\cdot(g-h)]}\le\eps$ for every $f\in\cF$; with $\cX=V\times V$ for a graph on $V$, $g$ its adjacency indicator, and $\cF$ the functions $\pm\ind_{S\times T}$, compare with the weak regularity lemma of \cref{ch:graphon}. (b) \textcite{TrevisanTulsianiVadhan2009} prove a $[0,1]$-valued, efficiently computable version, by LP duality as in Nisan's proof and by boosting, and derive from it the weak regularity lemma, a dense model theorem, and the hardcore lemma. Open-ended: find another structure-versus-randomness decomposition that is a min--max theorem in disguise; no general criterion is known. \hint{For (a), project $g$ onto the span of the functions chosen so far, in the inner product $\E[uv]$; while some $f\in\cF$ violates the bound, add it, and show that $\E[(g-h)^2]$ drops by more than $\eps^2$.}\label{ex:lpduality:regularity}
\end{exercises}

\notesandrefs

\Textcite{Kjeldsen2001} follows von Neumann's conception of the minimax theorem through its mathematical contexts, and \textcite{vonStengel2024} sorts out which of the duality and minimax theorems proves which, starting from the conversation recalled by \textcite{Dantzig1982}. \textcite{Schrijver1986} is the standard reference for polyhedra, total unimodularity, and LP algorithms.

\textcite{Schrijver2005} tells how the combinatorial zoo met linear programming. \textcite{FreundSchapire1996} develop the game-theoretic reading of boosting that \cref{ch:regret} turns into dynamics, and \textcite{TrevisanTulsianiVadhan2009} put regularity lemmas, dense model theorems, and the hardcore lemma under one min\nobreakdash--max roof.

\printchapterbib
\end{refsection}
